% 1999 Q15 井底抓斗提升模型图 (物理应用)
% 井身截面轮廓 (剖面阴影)
\draw[line width=1.1pt] (-2.2, 3.2) -- (-1.1, 3.2) -- (-1.1, -0.2) -- (0.8, -0.2) -- (0.8, 3.2) -- (1.6, 3.2);
\draw[line width=0.8pt] (-2.2, 2.8) -- (-1.4, 2.8) -- (-1.4, -0.5) -- (1.1, -0.5) -- (1.1, 2.8) -- (1.6, 2.8);
\draw[line width=0.8pt] (-2.2, 3.2) -- (-2.2, 2.8);
\draw[line width=0.8pt] (1.6, 3.2) -- (1.6, 2.8);

% 井口滑轮支架
\coordinate (P1) at (-1.7, 4.0);
\coordinate (P2) at (-0.2, 4.0);

% 滑轮 1
\draw[line width=0.9pt] (P1) circle (0.35);
\draw[line width=0.8pt] (-1.85, 3.2) -- (P1) -- (-1.55, 3.2);

% 滑轮 2
\draw[line width=0.9pt] (P2) circle (0.35);
\draw[line width=0.8pt] (-0.9, 3.2) -- (P2) -- (-0.5, 3.2);

% 缆绳
\draw[line width=1.0pt] (-1.7, 4.35) -- (-0.2, 4.35);
\draw[line width=1.0pt] (0.15, 4.0) -- (0.15, 1.2);

% 抓斗 (梯形)
\draw[line width=1.1pt] (-0.15, 1.2) -- (0.45, 1.2) -- (0.6, 0.4) -- (-0.3, 0.4) -- cycle;

% 右侧坐标轴 (高度 x, 0 到 30)
\draw[->, line width=0.9pt] (2.2, -0.5) -- (2.2, 4.2) node[above=1pt] {$x$};
\draw[line width=0.9pt] (2.0, -0.2) -- (2.4, -0.2) node[below=2pt] {$0$};
\draw[line width=0.9pt] (2.0, 3.2) -- (2.4, 3.2) node[right=3pt] {$30$};
